\section{讨论与局限}

\paragraph{候选级学习的解释。}
NELL23K 的官方 EM 为 91.52\%，高于原始 GRIP 论文的 87.74\%。同候选前 200 条问题上，本地 SC-CL 的 MRR 也高于 B1。这些观察与困难候选提供更清晰关系边界的设想相容。不过，前缀子集不能代表完整测试集；结构画像与基线违例度的各自作用还需要单信号消融和匹配训练预算的随机候选对照。

\paragraph{开放世界边界。}
train-known 过滤仅排除训练图中该实体对已有的关系。剩余候选可能在持出 split 或真实世界中成立，因此“结构反事实”是训练标签替换的操作名，不是经过验证的假事实。持出正例重叠与开放世界未判定需要分开统计；CoDEx 的实体替换硬负例也不能为本文的关系替换候选提供真值 \citep{safavi2020codex}。

\paragraph{成本与适用范围。}
冻结 bank 使候选来源可追溯，但需要离线计算 B1 的全关系分数，也不能随训练中的模型更新。现有资源记录只覆盖 NELL23K 的一次运行，尚不能判断收益是否抵消额外成本。方法针对固定实体对上的关系预测；图级组合问答需要重新定义反事实单位和评估目标。跨图适用性仍取决于 FB15K237-2、WN18RR 和 CoDEx-M 上的同契约实验。
